\documentclass[12pt]{article}
\usepackage[utf8]{inputenc}
\usepackage[spanish]{babel}
\usepackage{amsmath, amssymb}
\usepackage{geometry}
\geometry{margin=2.5cm}

\title{Modelo Examen Primer Parcial}
\author{}
\date{}

\begin{document}
\maketitle

\section*{Ejercicio 1}
Calcular el límite
\[
\lim_{x\to 0}\frac{\sqrt{1+\sen x}-e^{x/2}}{x^2}.
\]

\section*{Ejercicio 2}
Sea la función
\[
f(x)=
\begin{cases}
\dfrac{e^x-1}{x}, & x<0,\\[8pt]
a, & x=0,\\[8pt]
\dfrac{\sqrt{1+2x}-1}{x}+bx, & x>0,
\end{cases}
\qquad a,b\in\mathbb{R}.
\]
Determinar, de forma razonada, los valores de \(a\) y \(b\) para que la función sea:
\begin{enumerate}
    \item continua en \(x=0\),
    \item derivable en \(x=0\).
\end{enumerate}

\section*{Ejercicio 3}
Calcula \(\sqrt[5]{34}\) con un error menor que \(10^{-5}\) utilizando el desarrollo de Taylor adecuado.
\end{document}